\begin{abstract}
Dynamic-depth neural networks can trade accuracy for computation by selecting, per input, between a shallow (base) and a deep (super) execution path, but doing so requires a router that decides when the extra depth is worth its cost. 
Existing efficiency-aware routers are typically trained with a weighted sum of an accuracy term and a computation (FLOPs) term, frequently incorporating an additional uniformity or entropy term. 
This couples the operating point to training: the trade-off weight bakes a single accuracy–efficiency point into the model, while collapse onto one path must be prevented with auxiliary terms, and serving a new compute budget demands retraining. 
We reformulate depth routing as \emph{advantage regression}. 
Instead of learning a discrete policy, a lightweight head predicts, for each image, how much detection loss the deep path would save over the shallow path. 
Consequently, routing reduces to simply comparing this predicted advantage against a threshold at deployment time. 
A single trained router therefore produces an entire continuous accuracy–efficiency Pareto curve through a simple threshold sweep, with no trade-off hyperparameter and no collapse-prevention term. 
On KITTI (trained on detection images, evaluated on tracking videos), the proposed router dominates random, luminance, edge-density, and confidence-based baselines, and the result transfers to BDD100K.
The router adds negligible overhead, a few thousandths of a percent of the detector computation, and yields measurable on-device latency, energy, and frame-rate
improvements.
\end{abstract}

\begin{IEEEkeywords}
Autonomous Driving,
Adaptive inference, 
Conditional Computation, 
Dynamic Neural Networks, 
Efficient Deep Learning, 
Object Detection, 
Accuracy--Efficiency Trade-off, 
\end{IEEEkeywords}